\documentclass[a4paper,11pt]{uzreport}
\usepackage{booktabs}
\usepackage{tabularx}
\usetikzlibrary{arrows.meta,positioning}
\type{p}
\class{Sieci społecznościowe i systemy wieloagentowe}
\lab{Wieloagentowa giełda samochodowa}
\labnumber{Projekt}
\date{2026}
\supervisor{Dr inż. Marek Wróblewski}
\group{12INF-ISI-SD(L)}
\author{Olaf Bykowski\\\texttt{109217@stud.uz.zgora.pl}}
\hypersetup{hidelinks,pdftitle={Wieloagentowa giełda samochodowa}}
\pgfplotsset{compat=1.18}
\setcounter{secnumdepth}{2}
\renewcommand{\thesubsection}{\thesection.\arabic{subsection}}
\renewcommand{\thefigure}{\arabic{figure}}
\renewcommand{\thetable}{\arabic{table}}
\AtBeginDocument{\renewcommand{\thelstlisting}{\arabic{lstlisting}}}
\lstset{style=usreportlst,basicstyle=\ttfamily\footnotesize,
  columns=fullflexible,numbers=left,numberbychapter=false}
\setlength{\emergencystretch}{2em}
\newcommand{\code}[1]{\texttt{\detokenize{#1}}}
\newcommand{\testcase}[7]{%
  \subsubsection*{#1}
  \begin{description}[style=nextline,leftmargin=0pt,labelwidth=0pt,labelsep=0pt,font=\bfseries]
    \item[Warunki początkowe] #2
    \item[Agenci / komponenty] #3
    \item[Przebieg] #4
    \item[Oczekiwany rezultat] #5
    \item[Uzyskany rezultat] #6
    \item[Wynik] #7
  \end{description}}
\begin{document}
\maketitle
\tableofcontents
\clearpage

\section{Wprowadzenie}
\subsection{Cel i opis problemu}
Celem projektu jest zaprojektowanie i implementacja symulacji giełdy
samochodowej z wykorzystaniem Java, JADE/ACL, Elasticsearch oraz Kibany.
Niezależni agenci reprezentują sprzedających, kupujących i usługi doradcze.
Podstawowym problemem jest skoordynowanie autonomicznych decyzji zakupowych:
kilku kupujących może wybrać tę samą ofertę, lecz samochód może mieć tylko
jednego nabywcę. Jednocześnie należy respektować budżet, preferencje,
uzgodnioną cenę oraz ograniczenia czasowe rezerwacji i zwrotu.

\subsection{Zakres funkcjonalny}
System pobiera dane identyfikacyjne pojazdów z NHTSA vPIC, oczyszcza je
i uzupełnia brakujące parametry rynkowe w sposób deterministyczny.
Konfiguracja podstawowa obejmuje 10 sprzedających, po 8 samochodów,
5 kupujących, po 3 docelowe zakupy, oraz czterech agentów usługowych.
Łącznie w symulacji uczestniczy 19 agentów aplikacyjnych i 80 samochodów,
a docelowa liczba zachowanych zakupów wynosi 15.

Zakres obejmuje wyszukiwanie i porównywanie ofert, koszt całkowity,
kontrolę budżetu, konkurencję, rezerwacje, rekomendacje TOP~3,
negocjacje, reputację, analizę trendów, zwroty, wycenę rynkową
i analizę wielokryterialną (MCDA). Elasticsearch umożliwia zapis danych
i historii transakcji, a Kibana ich wizualizację.

\section{Architektura systemu}
\begin{figure}[htbp]
\centering
\resizebox{\textwidth}{!}{%
\begin{tikzpicture}[
  box/.style={draw,rounded corners,align=center,minimum height=10mm,
    text width=38mm,font=\small,fill=uzlogoblue!7},
  arr/.style={-{Latex},thick}]
  \node[box] (api) at (0,0) {NHTSA vPIC};
  \node[box] (import) at (5.5,0) {Importer\\i czyszczenie};
  \node[box] (json) at (11,0) {JSON\\80 samochodów};
  \node[box] (boot) at (11,-2.2) {Uruchomienie\\symulacji};
  \node[box] (agents) at (5.5,-2.2) {19 agentów\\aplikacyjnych};
  \node[box] (obs) at (0,-2.2) {SimulationObserver};
  \node[box] (es) at (0,-4.4) {Elasticsearch};
  \node[box] (kib) at (5.5,-4.4) {Kibana};
  \draw[arr] (api)--(import);
  \draw[arr] (import)--(json);
  \draw[arr] (json)--(boot);
  \draw[arr] (boot)--(agents);
  \draw[arr] (agents)-- node[above=1mm,font=\scriptsize,
    text width=14mm,align=center]{wyniki i\\zdarzenia} (obs);
  \draw[arr] (obs)-- node[right,font=\scriptsize]{MarketplaceEventSink} (es);
  \draw[arr] (es)--(kib);
\end{tikzpicture}}
\caption{Przepływ danych, uruchomienie agentów i zapis wyników.
Raportowanie do obserwatora jest niezależne od komunikacji ACL.}
\label{fig:arch}
\end{figure}

\begin{figure}[htbp]
\centering
\begin{tikzpicture}[
  box/.style={draw,rounded corners,align=center,minimum height=11mm,
    text width=36mm,font=\small,fill=uzlogoblue!7},
  both/.style={{Latex}-{Latex},thick},
  note/.style={font=\scriptsize,align=center,text width=36mm,
    above=2mm}]
  \node[box] (buyer) at (0,0) {5 $\times$ BuyerAgent};
  \node[box] (seller) at (10,4) {10 $\times$ SellerAgent\\SellerInventory};
  \node[box] (advisor) at (10,2) {AdvisorAgent};
  \node[box] (neg) at (10,0) {NegotiatorAgent};
  \node[box] (rep) at (10,-2) {ReputationAgent};
  \node[box] (trend) at (10,-4) {TrendAgent};
  \draw[both] (buyer.east)--(4,4)--
    node[note]{wyszukiwanie, rezerwacje,\\zakupy i zwroty} (seller.west);
  \draw[both] (buyer.east)--(4,2)--
    node[note]{rekomendacje TOP 3, MCDA\\i wycena} (advisor.west);
  \draw[both] (buyer.east)--
    node[note,pos=.73]{oferty cenowe\\i kontrpropozycje} (neg.west);
  \draw[both] (buyer.east)--(4,-2)--
    node[note]{oceny sprzedawców\\i zapytania o reputację} (rep.west);
  \draw[both] (buyer.east)--(4,-4)--
    node[note]{zdarzenia sprzedaży\\i wyniki analizy trendów} (trend.west);
\end{tikzpicture}
\caption{Komunikacja JADE/ACL. Strzałki dwukierunkowe oznaczają
żądania lub zdarzenia oraz odpowiedzi agentów.}
\label{fig:acl}
\end{figure}

Bootstrap symulacji realizuje \code{MarketplaceSimulationListener}.
Agenci usługowi rejestrują się w JADE Directory Facilitator (DF).
Kupujący odkrywa usługi dynamicznie i rozpoczyna pracę dopiero po
znalezieniu wszystkich wymaganych usług. Usługi doradcze nie przejmują
własności magazynu ani budżetu kupującego.

\code{SellerInventory} jest autorytatywnym właścicielem stanu pojazdów.
Operacje rezerwacji, zakupu, wygaśnięcia i zwrotu są synchronizowane.
Stan pojazdu przechodzi między \code{AVAILABLE}, \code{RESERVED}
i \code{SOLD}. \code{SimulationObserver} zbiera wyniki i przekazuje
transakcje do opcjonalnego \code{MarketplaceEventSink}.
Logika agentowa nie zależy od dostępności Elasticsearch.

\section{Opis agentów}
\subsection{SellerAgent --- sprzedający (10 instancji)}
\textbf{Funkcja:} zarządzanie ośmioma pojazdami i obsługa sprzedaży.
\textbf{Wejście:} kryteria wyszukiwania, identyfikator samochodu, kupujący,
uzgodniona cena i identyfikator transakcji.
\textbf{Wyjście:} najlepsza dostępna oferta, wynik rezerwacji, zakupu lub zwrotu.
\textbf{Komunikaty:} \code{car-search}, \code{car-reservation},
\code{car-purchase}, \code{car-return}.
\textbf{Decyzje:} wybór najtańszej zgodnej oferty lub odmowa; kontrola dostępności,
ceny minimalnej, właściciela rezerwacji i prawa do zwrotu.

\subsection{BuyerAgent --- kupujący (5 instancji)}
\textbf{Funkcja:} realizacja trzech zachowanych zakupów z ograniczonym budżetem.
\textbf{Wejście:} saldo, preferencje oraz odpowiedzi sprzedających i usług.
\textbf{Wyjście:} zapytania ACL, wybór oferty, płatność, ocena i zdarzenie sprzedaży.
\textbf{Komunikaty:} wszystkie ontologie opisane w sekcji~\ref{sec:acl}.
\textbf{Decyzje:} wybór rekomendacji mieszczącej się w budżecie,
rozpoczęcie negocjacji, ponowienie wyszukiwania po utracie oferty,
zakup i ewentualny zwrot. Po zwrocie kupuje pojazd zastępczy.

\subsection{AdvisorAgent --- doradca}
\textbf{Funkcja:} filtrowanie i wielokryterialna ocena ofert.
\textbf{Wejście:} oferty, kryteria kupującego i reputacje sprzedawców.
\textbf{Wyjście:} TOP~3 z wynikami MCDA i wycena pierwszej oferty.
\textbf{Komunikaty:} \code{car-recommendation}.
\textbf{Decyzje:} odrzucenie niezgodnych ofert, ranking według ważonej oceny;
przy remisie preferowana jest lepsza reputacja.

\subsection{NegotiatorAgent --- negocjator}
\textbf{Funkcja:} prowadzenie ograniczonej negocjacji ceny.
\textbf{Wejście:} oferta katalogowa, pierwsza propozycja kupującego i jego limit.
\textbf{Wyjście:} historia rund, kontrpropozycje i wynik negocjacji.
\textbf{Komunikaty:} \code{car-negotiation}.
\textbf{Decyzje:} akceptacja, kontrpropozycja lub odrzucenie; maksymalnie
trzy rundy i rabat do 15\%. Sprzedawca niezależnie weryfikuje uzgodnioną cenę.

\subsection{ReputationAgent --- reputacja}
\textbf{Funkcja:} agregacja ocen sprzedających.
\textbf{Wejście:} zapytanie o sprzedawców albo ocena 1--5 dla transakcji.
\textbf{Wyjście:} średnia, liczba ocen lub błąd walidacji.
\textbf{Komunikaty:} \code{reputation-query}, \code{seller-rating}.
\textbf{Decyzje:} przyjęcie poprawnej oceny potwierdzonej transakcji,
odrzucenie duplikatu lub nieuprawnionej oceny.

\subsection{TrendAgent --- analiza rynku}
\textbf{Funkcja:} analiza zachowanych sprzedaży.
\textbf{Wejście:} potwierdzone zdarzenia sprzedaży.
\textbf{Wyjście:} popularne nadwozie, model i paliwo, średnia cena i kierunek zmian.
\textbf{Komunikaty:} \code{market-sale-event}.
\textbf{Decyzje:} walidacja zdarzenia, aktualizacja agregatów i klasyfikacja trendu.
Trend krótkiej symulacji opisuje jej kolejne transakcje, a nie długookresowy
trend rzeczywistego rynku.

\section{Komunikacja JADE/ACL}\label{sec:acl}
Treść komunikatów jest kodowana jako JSON. Pole \code{ontology} określa
operację, a \code{conversationId}, \code{replyWith} i \code{inReplyTo}
pozwalają powiązać odpowiedź z odpowiednim żądaniem. Kupujący kontroluje
również nadawcę oraz termin odpowiedzi, aby nie przyjąć wiadomości z innej fazy.

\begin{table}[htbp]
\centering
\small
\begin{tabular}{p{.28\textwidth}p{.25\textwidth}p{.36\textwidth}}
\toprule Ontologia & Kierunek & Performatywy\\\midrule
\code{car-search} & Buyer $\to$ Seller & CFP; PROPOSE / REFUSE\\
\code{reputation-query} & Buyer $\to$ Reputation & QUERY\_REF; INFORM\\
\code{car-recommendation} & Buyer $\to$ Advisor & REQUEST; INFORM\\
\code{car-negotiation} & Buyer $\to$ Negotiator & REQUEST; INFORM\\
\code{car-reservation} & Buyer $\to$ Seller & REQUEST; AGREE / REFUSE\\
\code{car-purchase} & Buyer $\to$ Seller & REQUEST; INFORM / FAILURE\\
\code{seller-rating} & Buyer $\to$ Reputation & INFORM; INFORM / FAILURE\\
\code{market-sale-event} & Buyer $\to$ Trend & INFORM; INFORM / FAILURE\\
\code{car-return} & Buyer $\to$ Seller & REQUEST; INFORM / REFUSE\\
\bottomrule
\end{tabular}
\caption{Ontologie, kierunki żądań i performatywy odpowiedzi.}
\end{table}

Poniższy zapis przedstawia \textbf{przykładową parę komunikatów},
ilustrującą rezerwację; nie jest zrzutem konkretnej rozmowy z logu.
\begin{lstlisting}[caption={Żądanie rezerwacji ACL i treść JSON}]
(REQUEST
 :sender Buyer01
 :receiver (set Seller08)
 :ontology car-reservation
 :conversation-id reservation-example-01
 :reply-with reservation-request-01
 :content "{\"carId\":\"CAR-062\",\"buyerName\":\"Buyer01\",\"agreedTotalPln\":124590}")
\end{lstlisting}
\begin{lstlisting}[caption={Odpowiedź powiązana z żądaniem}]
(AGREE
 :sender Seller08
 :receiver (set Buyer01)
 :ontology car-reservation
 :conversation-id reservation-example-01
 :in-reply-to reservation-request-01
 :content "<JSON wyniku rezerwacji>")
\end{lstlisting}
Zawartość pola odpowiedzi oznaczono schematycznie. Sprzedający nie ufa
samej treści JSON: porównuje \code{buyerName} z tożsamością nadawcy ACL,
sprawdza dostępność i cenę minimalną. Nieudana rezerwacja nie oznacza
zakończenia pracy kupującego, lecz konieczność wyboru innej oferty.

\section{Najważniejsze elementy implementacji}
\subsection{Pozyskiwanie i model danych}
Importer odrzuca niekompletne rekordy oraz duplikaty. Marka, model, rok
i kategoria pojazdu pochodzą z NHTSA vPIC. Cena, koszt dodatkowy,
silnik, pojemność, przebieg, moc, spalanie, skrzynia i napęd są uzupełniane
deterministycznie, ponieważ vPIC nie jest źródłem ofert handlowych.
Zbiór 80 rekordów zapisywany jest jako \code{data/cars.json}.
Oznacza to, że ceny w projekcie są danymi symulacyjnymi, a nie
zaobserwowanymi cenami rzeczywistych ogłoszeń.

Model zawiera identyfikator, markę, model, nadwozie, typ silnika,
pojemność, rok, cenę i koszt dodatkowy, a także rozszerzone parametry
techniczne. Koszt całkowity jest wspólną podstawą filtrowania,
porównywania i obciążania budżetu:
\[
K_{\text{całk.}}=P_{\text{samochodu}}+K_{\text{dodatkowy}}.
\]
Przykładowo cena 98\,000 PLN i opłaty 5\,000 PLN dają 103\,000 PLN,
zatem kupujący z budżetem 100\,000 PLN nie może dokonać takiego zakupu.

\subsection{Rezerwacja, współbieżność i zakup}
Rezerwacja domyślnie trwa 5 sekund. Jej właściciel uzyskuje prawo
zakupu po uzgodnionej cenie; pozostali kupujący otrzymują odmowę.
Po upływie czasu bez zakupu samochód wraca do \code{AVAILABLE}.
Synchronizacja obejmuje całą operację sprawdzenia warunków i zmiany stanu,
a nie tylko osobny odczyt dostępności. Eliminuje to wyścig prowadzący
do podwójnej rezerwacji lub sprzedaży.

\subsection{Wybrane fragmenty kodu}
Poniższe listingi są fragmentami rzeczywistej implementacji. Pominięto
kod otaczający, aby wyeksponować najważniejsze mechanizmy.

\begin{lstlisting}[language=Java,caption={BuyerAgent.java: utworzenie żądania rezerwacji (wiersze 397--408)}]
private void sendReservation() {
    ACLMessage request = new ACLMessage(ACLMessage.REQUEST);
    request.addReceiver(selectedSeller);
    request.setOntology(MarketplaceProtocol.RESERVATION);
    request.setConversationId(conversationId);
    expectedReplyToken = UUID.randomUUID().toString();
    request.setReplyWith(expectedReplyToken);
    request.setContent(codec.encode(new MarketplaceProtocol.ReservationRequest(
            selectedOffer.car().id(), getLocalName(), agreedTotalPln)));
    send(request);
    state = WAITING_FOR_RESERVATION;
}
\end{lstlisting}

\begin{lstlisting}[language=Java,caption={SellerInventory.java: atomowa kontrola i zapis rezerwacji (wiersze 45--67)}]
public synchronized ReservationResult reserve(
        String carId,
        String buyerName,
        long agreedTotalPln,
        long durationMillis) {
    releaseExpiredReservations();
    InventoryItem item = items.get(carId);
    if (item == null) {
        return new ReservationResult(false, carId, "UNKNOWN_CAR");
    }
    if (item.status != CarStatus.AVAILABLE) {
        return new ReservationResult(false, carId, "CAR_UNAVAILABLE");
    }
    if (agreedTotalPln < NegotiationService.minimumTotal(item.car)
            || agreedTotalPln > item.car.totalCostPln()) {
        return new ReservationResult(false, carId, "INVALID_NEGOTIATED_PRICE");
    }
    item.status = CarStatus.RESERVED;
    item.reservedBy = buyerName;
    item.reservedTotalCostPln = agreedTotalPln;
    item.reservedUntilEpochMillis = clock.millis() + durationMillis;
    return new ReservationResult(true, carId, "RESERVED");
}
\end{lstlisting}

\begin{lstlisting}[language=Java,caption={SellerInventory.java: kontrola właściciela przed sprzedażą (wiersze 69--76)}]
public synchronized Optional<Transaction> purchase(String carId, String buyerName) {
    releaseExpiredReservations();
    InventoryItem item = items.get(carId);
    if (item == null || item.status != CarStatus.RESERVED || !buyerName.equals(item.reservedBy)) {
        return Optional.empty();
    }
    item.status = CarStatus.SOLD;
    long transactionTotalCostPln = item.reservedTotalCostPln;
\end{lstlisting}

\begin{lstlisting}[language=Java,caption={RecommendationService.java: normalizacja kryteriów kosztowych i korzyści (wiersze 79--85)}]
private double lowerIsBetter(double value, Range range) {
    return range.isConstant() ? 1.0 : (range.maximum() - value) / (range.maximum() - range.minimum());
}

private double higherIsBetter(double value, Range range) {
    return range.isConstant() ? 1.0 : (value - range.minimum()) / (range.maximum() - range.minimum());
}
\end{lstlisting}
Stała wartość kryterium daje wszystkim ofertom ocenę 1, dzięki czemu nie
dochodzi do dzielenia przez zero. Normalizacja jest względna wobec
badanego zbioru; wyników MCDA z różnych zbiorów nie należy traktować
jako bezwzględnej miary jakości pojazdu.


\subsection{Rekomendacje i MCDA}
\code{RecommendationService} normalizuje wartości w bieżącym zbiorze
zgodnych ofert, a następnie oblicza ważoną sumę ocen cząstkowych:
\[
 S(a)=0{,}30s_{\mathrm{cena}}+
 0{,}20s_{\mathrm{rok}}+
 0{,}15s_{\mathrm{koszt}}+
 0{,}10s_{\mathrm{paliwo}}+
 0{,}10s_{\text{pojemność}}+
 0{,}15s_{\mathrm{spalanie}}.
\]
Wagi sumują się do 100\%. Cena, koszt całkowity i spalanie są kryteriami
kosztowymi, natomiast nowszy rok i większa pojemność zwiększają ocenę.
Ocena paliwa uwzględnia jawne kryterium kupującego, a przy jego braku
stosuje kolejność ELECTRIC, HYBRID, PETROL, DIESEL.
Wynik mieści się w przedziale $[0,1]$,
a reputacja rozstrzyga remisy. Typ nadwozia jest uwzględniany przy
filtrowaniu ofert. Ranking TOP~3 jest faktycznie wykorzystywany przy wyborze
oferty do negocjacji i zakupu.

\subsection{Negocjacje, reputacja i wycena}
Negocjator ogranicza rozmowę do trzech rund. Dopuszczalny rabat wynosi
maksymalnie 15\% ceny samochodu; koszt dodatkowy pozostaje częścią
całkowitego kosztu transakcji. Reputacja jest średnią ocen 1--5,
z kontrolą historii transakcji i odrzucaniem ponownych ocen.

Wycena wykorzystuje porównywalne pojazdy o takim samym nadwoziu i paliwie
oraz roku różniącym się najwyżej o dwa lata. Porównanie ceny ze średnią
pozwala zaklasyfikować odchylenie do przedziałów 0--5\%, 5--15\%
i powyżej 15\%. Jest to metoda porównawcza na zbiorze symulacyjnym,
nie wytrenowany model predykcyjny rzeczywistej wartości pojazdu.

\subsection{Zwroty i historia}
Zwrot może zgłosić wyłącznie kupujący zapisany w transakcji i tylko
w czasie określonym przez \code{simulation.return-window-ms}.
Po akceptacji samochód wraca do \code{AVAILABLE}, kupujący otrzymuje
zwrot środków, a historia transakcji uzyskuje status \code{RETURNED}.
Scenariusz demonstracyjny obejmuje zakup zamiennika, dlatego po zwrocie
nadal pozostaje 15 zachowanych zakupów.

\section{Testy}
\subsection{Metodyka i źródło wyników}
Poniższe scenariusze przeniesiono z \code{docs/SCENARIUSZE_TESTOWE.md}.
Dokument źródłowy podaje wyniki automatyczne uzyskane poleceniem
\code{./gradlew test} i wyniki integracyjne z pełnej symulacji opisanej
w \code{docs/DEMO.md}. Testy reguł biznesowych sprawdzają komponenty
używane przez agentów; nie każdy z nich uruchamia platformę JADE.
Wyniki PASS poniżej są wynikami zapisanymi w dokumentacji projektu,
a nie deklaracją ponownego uruchomienia testów podczas składu sprawozdania.
Uwzględniono wszystkie 17 scenariuszy, co przekracza minimum 10
wskazane w punktacji końcowej.

\testcase{T01 --- Wyszukanie pasującego auta}
{Dostępne jest auto zgodne z kryteriami.}
{BuyerAgent, SellerAgent.}
{Kupujący wysyła \code{car-search}.}
{Sprzedający zwraca PROPOSE z najtańszą zgodną ofertą.}
{Oferta została zwrócona i użyta w rankingu.}{PASS}

\testcase{T02 --- Brak pasującego auta}
{Żadne dostępne auto nie spełnia filtrów.}
{BuyerAgent, SellerAgent.}
{Kupujący wysyła zapytanie z niemożliwymi kryteriami.}
{Odpowiedź REFUSE / NO\_OFFER.}
{Brak oferty obsłużono bez przerwania symulacji.}{PASS}

\testcase{T03 --- Zakup w ramach budżetu}
{Saldo przekracza uzgodniony koszt całkowity.}
{BuyerAgent, SellerAgent.}
{Rezerwacja, zakup i płatność.}
{Auto otrzymuje SOLD, a saldo maleje o koszt transakcji.}
{Transakcja została utworzona; saldo jest poprawne.}{PASS}

\testcase{T04 --- Zakup przekraczający budżet}
{Koszt całkowity przekracza saldo.}
{\code{BuyerAccount}, komponent budżetu kupującego.}
{Próba płatności ponad dostępne środki.}
{Odrzucenie operacji bez zmiany salda.}
{\code{IllegalStateException}; saldo pozostało bez zmian.}{PASS}

\testcase{T05 --- Konkurencja o jeden samochód}
{Kilku kupujących wybiera tę samą ofertę.}
{Kilka BuyerAgent i jeden SellerAgent.}
{Równoległe żądania rezerwacji.}
{Tylko jedno AGREE; pozostałe żądania otrzymują REFUSE.}
{W logu odnotowano odrzucone konkurencyjne rezerwacje.}{PASS}

\testcase{T06 --- Rezerwacja samochodu}
{Auto AVAILABLE, poprawna uzgodniona cena.}
{BuyerAgent, SellerAgent.}
{Wysłanie \code{car-reservation}.}
{Status RESERVED i przypisanie właściciela rezerwacji.}
{Rezerwacja została zaakceptowana.}{PASS}

\testcase{T07 --- Wygaśnięcie rezerwacji}
{Auto zarezerwowane bez późniejszego zakupu.}
{\code{SellerInventory}, komponent sprzedającego.}
{Przesunięcie zegara poza pięciosekundowy okres rezerwacji.}
{Automatyczny powrót do AVAILABLE.}
{Test sterowanego zegara potwierdził zwolnienie pojazdu.}{PASS}

\testcase{T08 --- Udana negocjacja}
{Możliwe porozumienie w granicach maksymalnego rabatu 15\%.}
{BuyerAgent, NegotiatorAgent; test reguł \code{NegotiationService}.}
{Oferta kupującego i kontrpropozycja.}
{Akceptacja w najwyżej trzech rundach.}
{Cena została zaakceptowana i użyta w transakcji. Test usługi
potwierdza kontrpropozycję 133\,500 PLN przy pierwszej ofercie
119\,000 PLN i limicie 148\,000 PLN.}{PASS}

\testcase{T09 --- Odrzucona negocjacja}
{Limit kupującego jest niższy niż minimalna dopuszczalna cena.}
{BuyerAgent, NegotiatorAgent; \code{NegotiationService}.}
{Negocjacja z pierwszą ofertą 110\,000 PLN i limitem 120\,000 PLN.}
{Odrzucenie najpóźniej po trzech rundach.}
{\code{accepted=false}; trzy rundy i brak uzgodnionej ceny.}{PASS}

\testcase{T10 --- Rekomendacja TOP 3}
{Co najmniej cztery oferty SUV spełniają kryteria.}
{BuyerAgent, AdvisorAgent; \code{RecommendationService}.}
{Przekazanie ofert i preferencji doradcy.}
{Trzy oferty o najwyższej ocenie MCDA w poprawnej kolejności.}
{Zwrócono CAR-001, CAR-002 i CAR-003; niezgodny SEDAN odrzucono.}{PASS}

\testcase{T11 --- Ocena sprzedawcy}
{Potwierdzone transakcje umożliwiają wystawienie ocen.}
{BuyerAgent, ReputationAgent; \code{ReputationService}.}
{Oceny 5 i 4 dla sprzedawcy oraz ponowna ocena pierwszej transakcji.}
{Średnia 4,5, dwie oceny i odrzucenie duplikatu.}
{Średnia oraz licznik są poprawne; duplikat odrzucono.}{PASS}

\testcase{T12 --- Analiza trendu}
{Cztery sprzedaże o rosnących cenach.}
{BuyerAgent, TrendAgent; \code{TrendService}.}
{Kolejne zdarzenia sprzedaży: 100\,000, 105\,000, 130\,000 i 140\,000 PLN.}
{Poprawne popularne kategorie, średnia i trend RISING.}
{Średnia wyniosła 118\,750 PLN; najpopularniejsze są SUV,
Toyota RAV4 i PETROL; trend RISING.}{PASS}

\testcase{T13 --- Zwrot w terminie}
{Od zakupu nie upłynęło pięć sekund.}
{BuyerAgent, SellerAgent.}
{Właściciel wysyła \code{car-return}.}
{Zwrot środków, AVAILABLE i historia RETURNED.}
{Wykonano pełny zwrot; kupujący następnie kupił zamiennik.}{PASS}

\testcase{T14 --- Zwrot po terminie}
{Od zakupu minęło ponad pięć sekund.}
{BuyerAgent, SellerAgent.}
{Żądanie zwrotu po upływie okna czasowego.}
{Odmowa i pozostawienie auta w stanie SOLD.}
{Otrzymano \code{RETURN_WINDOW_EXPIRED}.}{PASS}

\testcase{T15 --- Określenie wartości rynkowej}
{Oferta oraz dwa porównywalne pojazdy.}
{AdvisorAgent; \code{ValuationService}.}
{Obliczenie średniej i procentowego odchylenia ceny.}
{Poprawna kwalifikacja do przedziałów 0--5, 5--15 i powyżej 15\%.}
{Zweryfikowano granice przedziałów klasyfikacji.}{PASS}

\testcase{T16 --- Wybór metodą MCDA}
{Oferty różnią się ceną, rokiem, paliwem, pojemnością i spalaniem.}
{BuyerAgent, AdvisorAgent; \code{RecommendationService}.}
{Normalizacja cech i obliczenie ważonej sumy.}
{Ranking zgodny z wagami 30/20/15/10/10/15.}
{Kolejność rekomendacji i zakres ocen 0--1 są poprawne.}{PASS}

\testcase{T17 --- Jednoczesna komunikacja wielu agentów}
{10 sprzedających, 5 kupujących i 80 samochodów.}
{Wszystkie agenty aplikacyjne.}
{Uruchomienie pełnej symulacji.}
{15 zachowanych zakupów, bez podwójnej sprzedaży i pozostałych rezerwacji.}
{15 SOLD, 0 RESERVED; każdy kupujący ma trzy auta.}{PASS}

\subsection{Podsumowanie testów}
Dokumentacja źródłowa podaje 17 scenariuszy zaliczonych i zero niezaliczonych.
Scenariusze obejmują zarówno ścieżki poprawne, jak i odmowę z powodu
braku ofert, niewystarczającego budżetu, konkurencji, niedopuszczalnej
ceny lub przekroczenia czasu zwrotu.


\section{Wyniki}
\subsection{Transakcje i stan końcowy}
Zgodnie z udokumentowanym przebiegiem podstawowym każdy z pięciu kupujących
zachował trzy samochody. Końcowy stan magazynów przedstawiono
w tabeli~\ref{tab:runs}.
\begin{table}[htbp]
\centering
\begin{tabular}{lrr}
\toprule
Wielkość & Przebieg podstawowy & Przebieg ze zwrotem\\
\midrule
Samochody na wejściu & 80 & 80\\
Rekordy zakupu & 15 & 16\\
Zwroty & 0 & 1\\
Zachowane zakupy & 15 & 15\\
AVAILABLE & 65 & 65\\
RESERVED & 0 & 0\\
SOLD & 15 & 15\\
\bottomrule
\end{tabular}
\caption{Wyniki przebiegów opisanych w dokumentacji projektu.}
\label{tab:runs}
\end{table}

Przebieg demonstracyjny uruchamiany z
\code{--simulation.demonstrate-return=true} utworzył 16 rekordów zakupu
i jeden zwrot. Zakup zastępczy przywrócił docelowy stan.
Każdy przebieg zapisu do Elasticsearch ma osobny \code{runId},
co zapobiega nadpisywaniu historii kolejnych symulacji.

\subsection{Rekomendacje i wyniki MCDA}
Repozytorium nie zawiera zapisanego logu pełnej symulacji z liczbowymi
wynikami poszczególnych ofert. Dlatego poniższe szczegółowe dane
pochodzą z deterministycznych przypadków testowych, a nie z przebiegu
80 samochodów. Wyniki całej symulacji podano osobno w tabeli~\ref{tab:runs}.

W \code{RecommendationServiceTests} cztery zgodne oferty SUV kosztują
łącznie 120\,000, 130\,000, 140\,000 i 150\,000 PLN.
Każda ma koszt dodatkowy 2\,000 PLN, rok 2024, pojemność 1998 cm$^3$,
paliwo PETROL i spalanie 7,0 l/100 km. Oferta SEDAN za 90\,000 PLN
jest odrzucana mimo niższej ceny.

\begin{table}[htbp]
\centering
\begin{tabular}{lrrl}
\toprule
Oferta & Koszt całkowity [PLN] & Wynik MCDA & Decyzja\\
\midrule
CAR-001 & 120\,000 & 1,00 & TOP 1\\
CAR-002 & 130\,000 & 0,85 & TOP 2\\
CAR-003 & 140\,000 & 0,70 & TOP 3\\
CAR-004 & 150\,000 & 0,55 & Poza TOP 3\\
CAR-005 & 90\,000 & --- & Niezgodne nadwozie\\
\bottomrule
\end{tabular}
\caption{Ranking testowego zbioru ofert. Oceny liczbowe wyliczono
z danych testu i wzoru implementacji; asercje testu sprawdzają
kolejność TOP~3 oraz zakres ocen, nie te konkretne wartości.}
\end{table}

Cena i koszt całkowity odpowiadają tu łącznie za 45\% oceny.
Pozostałe cechy są jednakowe, więc każda oferta uzyskuje za nie 0,55.
Dla CAR-002 normalizacja ceny oraz kosztu daje $2/3$, stąd
$S=0{,}55+0{,}45\cdot2/3=0{,}85$.
Test remisu przy jednakowych parametrach preferuje Seller02
z reputacją 5,0 przed Seller01 z reputacją 3,0.

\subsection{Negocjacje}
W \code{NegotiationServiceTests} samochód ma cenę 145\,000 PLN
i koszt dodatkowy 3\,000 PLN. Minimalny koszt po rabacie wynosi
$145\,000\cdot0{,}85+3\,000=126\,250$ PLN.
Oferta kupującego 119\,000 PLN przy limicie 148\,000 PLN prowadzi
do kontrpropozycji i akceptacji 133\,500 PLN już w pierwszej rundzie.
Oszczędność względem katalogowego kosztu 148\,000 PLN wynosi
14\,500 PLN. Drugi przypadek, z limitem 120\,000 PLN,
kończy się odrzuceniem po trzech rundach.

\subsection{Reputacja i trendy}
Test reputacji zapisuje dla Seller01 oceny 5 i 4 z dwóch transakcji.
Uzyskana średnia to $(5+4)/2=4{,}5$, a licznik ocen wynosi 2.
Ponowna ocena tej samej transakcji jest odrzucana.

Test trendów przetwarza cztery transakcje o wartościach 100\,000,
105\,000, 130\,000 i 140\,000 PLN. Wynik obejmuje cztery sprzedaże,
średnią 118\,750 PLN, najpopularniejsze nadwozie SUV, model Toyota RAV4,
paliwo PETROL i kierunek RISING. Są to agregaty kontrolowanego
zbioru testowego, a nie estymacja trendu rzeczywistego rynku.

\subsection{Elasticsearch i Kibana}
Plik \code{kibana/dashboard.ndjson} zawiera dwa widoki danych, pięć
wizualizacji Vega i dashboard \emph{SSISW --- giełda samochodowa}.
Wizualizacje przedstawiają liczbę pojazdów według marki, nadwozia
i paliwa, rozkład roku produkcji oraz transakcje według statusu
COMPLETED / RETURNED. Odwołują się do indeksów \code{cars}
i \code{transactions}; filtr \code{runId} pozwala wybrać przebieg.
Zrzuty działającego dashboardu nie zostały dostarczone z materiałami,
więc ich nie zastępowano przykładowymi obrazami.


\subsection{Interpretacja i ograniczenia}
Wszystkie zachowane transakcje w udokumentowanym przebiegu otrzymały ocenę
i zostały przekazane do analizy trendów. Testy potwierdzają powiązanie
rankingu MCDA z wyborem oferty, wykorzystanie uzgodnionej ceny przy
transakcji, aktualizację reputacji i poprawną obsługę zwrotu.
Wyniki krótkiej symulacji nie stanowią reprezentatywnego badania rynku
samochodowego. Kolejność działań współbieżnych agentów może się zmieniać
między uruchomieniami, nawet gdy zbiór danych wejściowych jest ten sam.

\section{Wnioski}
Projekt obejmuje wszystkie 21 minimalnych wymagań funkcjonalnych
wskazanych w zadaniu: wymaganą liczbę agentów i pojazdów, komunikację
ACL, model i koszty, budżet, konkurencję bez podwójnej sprzedaży,
rezerwacje, doradcę i TOP~3, negocjatora, reputację, trendy,
zwroty, wycenę, MCDA oraz scenariusze testowe.

Najważniejszym problemem była jednoczesna rywalizacja kupujących
o ten sam pojazd. Rozwiązaniem jest synchronizowany, autorytatywny
magazyn sprzedającego oraz rezerwacja z ograniczeniem czasu.
Integralność komunikacji wzmacniają identyfikatory rozmów, terminy
odpowiedzi, walidacja nadawcy i kontrola historii transakcji.
Rozdzielenie logiki agentowej od zapisu zdarzeń pozwala prowadzić
symulację również bez Elasticsearch.

Możliwe kierunki rozwoju obejmują trwałe przechowywanie reputacji,
wycenę opartą na rzeczywistych ofertach i modelu regresyjnym,
zróżnicowanie strategii kupujących, rozproszenie kontenerów JADE
oraz analizę trendów w dłuższych okresach. Przed oddaniem sprawozdania
należy potwierdzić dane strony tytułowej i dołączyć własne zrzuty
wizualizacji Kibany, jeżeli mają stanowić dowód prezentacji środowiska.

\begin{thebibliography}{9}
\bibitem{requirements} \emph{Sieci społecznościowe i systemy wieloagentowe,
Laboratorium -- projekt aplikacji wieloagentowej}, 2026,
w szczególności wytyczne Lab 14 i kryteria oceny.
\bibitem{project} Dokumentacja projektu:
\code{docs/SPRAWOZDANIE.md}, \code{docs/SCENARIUSZE_TESTOWE.md},
\code{docs/DEMO.md} oraz kod źródłowy dołączony do projektu.
\bibitem{vpic} NHTSA, Vehicle Product Information Catalog (vPIC),
\url{https://vpic.nhtsa.dot.gov/}.
\end{thebibliography}
\end{document}
